\documentclass[11pt,a4paper]{article}
  \usepackage[T1]{fontenc}
  \usepackage[utf8]{inputenc}
  \usepackage{lmodern}
  \usepackage[margin=2.2cm]{geometry}
  \usepackage{amsmath,amssymb,amsthm}
  \usepackage{booktabs}
  \usepackage{longtable}
  \usepackage{array}
  \usepackage{xcolor}
  \usepackage{enumitem}
  \usepackage{fancyhdr}
  \usepackage{titlesec}
  \usepackage[hidelinks]{hyperref}
  \usepackage{caption}
  \usepackage{tabularx}
  \usepackage{multirow}

  \definecolor{rulecol}{RGB}{180,30,40}
  \definecolor{softbg}{RGB}{250,244,244}

  \newcommand{\reviewer}[1]{\vspace{1.0em}\noindent\textbf{\color{rulecol}#1}\par\vspace{0.25em}}
  \newcommand{\ourans}{\vspace{0.35em}\noindent\textbf{Response.}\ }
  \newcommand{\change}[1]{\vspace{0.25em}\noindent\textbf{Change in the manuscript.} #1\par\vspace{0.2em}}
  \newcommand{\loc}[1]{\textbf{Location:} #1\par}
  \usepackage{needspace}
  \setlength{\emergencystretch}{3em}

  % Point n: the reviewer's comment in full, then one sub-point per request,
  % each with the question it answers.
  \newcommand{\point}[3]{\par\vspace{1.2em}\Needspace{8\baselineskip}%
    \noindent{\large\bfseries\color{rulecol}Point #1: #2}\par\nopagebreak
    \noindent{\small\itshape Reviewer~3, comment #3}\par\nopagebreak}
  \newenvironment{reviewercomment}{\begin{quote}\small\itshape}{\end{quote}}
  \newcommand{\subpoint}[3]{\par\vspace{0.9em}\Needspace{6\baselineskip}%
    \noindent{\bfseries Point #1: #2}\par\nopagebreak\vspace{0.15em}
    \noindent\textbf{Question.}\ \emph{#3}\par\nopagebreak\vspace{0.2em}
    \noindent\textbf{Response.}\ }

  \title{\vspace{-2.2cm}\textbf{Response to the Editor and Reviewers}\\[0.35em]
  \large Revised manuscript \emph{``Conditional Generative Modelling of Thai Stock-Return Forecasts for Economic Applications''}\\[0.6em]
  \normalsize Response letter, sixth version: the leakage-free re-run, point by point}
  \author{}
  \date{}

  \pagestyle{fancy}
  \fancyhf{}
  \fancyhead[L]{Response to the Editor and Reviewers}
  \fancyhead[R]{Revised manuscript}
  \fancyfoot[C]{\thepage}
  \renewcommand{\headrulewidth}{0.4pt}
  \setlength{\parskip}{0.45em}
  \setlength{\parindent}{0pt}

\begin{document}
\maketitle

\section*{0. Summary of this letter}

This letter answers reviewer comments R3.M2 (data leakage in the public
implementation) and R3.M3 (market-regime definitions and evaluation
periods), and R3.M4 (the Diebold--Mariano statistics). It reports the
re-run that was carried out for
the revision: every model was retrained with strictly chronological splits,
with the normalisation fitted on the training part only, with early stopping
driven by a validation block, and with a final test block that was read
once, after the model was frozen.

The letter is organised point by point. Each point restates the
reviewer's comment in full, and each sub-point (1.1, 1.2, \dots) quotes the
request it answers before the response.

\begin{center}\small
\begin{tabular}{llrrl}
\toprule
Regime & Window & Stocks & Test dates & Observations per stock \\
\midrule
crash & 5 & 39 & 2020-02-17--2020-03-30 & 31--31 \\
crash & 10 & 39 & 2020-02-17--2020-03-30 & 31--31 \\
crash & 30 & 39 & 2020-02-17--2020-03-30 & 31--31 \\
normal & 5 & 39 & 2023-11-08--2024-04-03 & 101--101 \\
normal & 10 & 39 & 2023-11-08--2024-04-03 & 101--101 \\
normal & 30 & 39 & 2023-11-08--2024-04-03 & 101--101 \\
\bottomrule
\end{tabular}
\end{center}

The number of forecast observations is now read from the data instead of
assumed, and it is identical across windows within a regime, so the
comparisons that the manuscript draws across look-back windows are made
on the same dates.

\section*{1. Response to Reviewer 3, comment R3.M2}

\point{1}{Data leakage in the public implementation}{R3.M2}
\begin{reviewercomment}
Potential data leakage in the publicly available implementation. The repository linked in the manuscript raises a serious reproducibility concern. Its publicly accessible README explicitly notes that, in \texttt{training/catBoost\_conGAN.py}, a \texttt{MinMaxScaler} was fitted to all sliding-window observations before the training/test split, thereby incorporating the price range of the test period into preprocessing. If this procedure was used to generate the submitted results, it would contaminate the out-of-sample evaluation with future information.

The repository currently describes a related, differently titled retrospective audit, so the correspondence between those files and the present manuscript should not be assumed without clarification. Please identify the exact repository commit and scripts underlying each reported table and figure, explain when normalization parameters are fitted, and demonstrate that hyperparameter selection and early stopping use training/validation information only. I strongly recommend rerunning the experiments using strictly chronological splits and training-only preprocessing, with an untouched final test set and an auditable replication package. The relevant repository is \url{https://github.com/kkabinsky/conGANs_SET50}.
\end{reviewercomment}
\ourans
We accept the concern and we have re-run the experiments. The three
defects the reviewer identified in the archived implementation are
corrected, and every number in the revised empirical section now comes
from the corrected run.

\subpoint{1.1}{What the archived implementation did, and what it is now}{The repository's README notes that, in \texttt{training/catBoost\_conGAN.py}, a \texttt{MinMaxScaler} was fitted to all sliding-window observations before the training/test split, thereby incorporating the price range of the test period into preprocessing. If this procedure was used to generate the submitted results, it would contaminate the out-of-sample evaluation with future information.}
The table compares the archived implementation with the re-run, item by item.

\begin{center}\small
\begin{tabular}{p{0.30\textwidth}p{0.30\textwidth}p{0.30\textwidth}}
\toprule
Item & Archived implementation & This re-run \\
\midrule
Normalisation & \texttt{MinMaxScaler.fit\_transform} on every sliding window, before the split (lines 153--160) & the scaler is fitted on the training windows only; the validation and test blocks are transformed with those constants \\
Split & index offsets \texttt{train = 905 + kk}, \texttt{test = next 100} (lines 165--170), no validation block & a chronological split stated as dates: training, then a validation block, then the final test block \\
Early stopping & \texttt{eval\_set=(X\_test, y\_test\_orig)} (lines 234--240) & \texttt{eval\_set=(X\_val, y\_val)}; the test block is read once, at the end \\
Generator selection & the last epoch & the epoch with the best validation loss \\
Stored forecasts & per stock only for CatBoost, the cGAN and the hybrids & every model, per stock, per regime, per window \\
\bottomrule
\end{tabular}
\end{center}

\subpoint{1.2}{When the normalisation parameters are fitted}{Please explain when the normalization parameters are fitted.}
In the archived script the two \texttt{MinMaxScaler} objects were fitted on
the flattened window matrix \emph{before} the train/test split, so the price
range of the test period entered the training data. In the re-run the scaler
is fitted on the training windows only, and the constants actually used are
recorded for every stock, so the difference is inspectable and not merely
asserted.

\subpoint{1.3}{Evidence that selection uses training and validation only}{Please demonstrate that hyperparameter selection and early stopping use training/validation information only.}
Hyperparameters are the ones of the archived run (CatBoost: 10000 iterations,
learning rate 0.01, depth 6, \texttt{early\_stopping\_rounds}=50,
\texttt{use\_best\_model}=True). The only change is the set the early
stopping watches: the validation block. The generator is likewise selected by
validation loss. The test block is touched once, when the frozen model
forecasts it.
The cGAN is trained for 100 epochs (Adam, learning rate $2\times10^{-4}$,
batch size 64), and after every epoch the generator is scored on the
validation block; the generator with the lowest validation RMSE is kept.
Over the 234 cells the kept generator comes from epoch 0 to 99 (median 23); in 40 cells it is the generator after the first epoch, and in 94 cells it comes from the first five epochs. The per-cell record is in \nolinkurl{data/training_log/} of the release.

\subpoint{1.4}{The scripts and their hashes}{Please identify the exact repository commit and scripts underlying each reported table and figure. The repository currently describes a related, differently titled retrospective audit, so the correspondence between those files and the present manuscript should not be assumed without clarification.}
The submitted results were produced by per-stock copies of one training
script, \texttt{new\_cat\_conGAN\_modified.py}, kept in \texttt{D:\textbackslash CatBoost}: 39 copies for $W = 5$ (\texttt{window5\textbackslash SET50}), 39 for $W = 10$ (\texttt{window10\textbackslash SET50}) and 39 for $W = 30$ (\texttt{SET50}), one folder per stock.
Once the three lines that name the input file, set the window size and set
the number of iterations are removed, all 117 copies are
identical, so the three defects of Point~1.1 are at the same lines in every
copy.
The DAX version is \texttt{catBoost\_conGAN.py}.
One copy does not match its folder: the script of SAWAD in the $W = 5$ folder sets \texttt{window\_size = 30}.

None of these scripts was ever committed. The local repository that holds
them (remote \texttt{github.com/svmlstmganproject/CatBoost.git}) is at commit \texttt{06bba6720374} of 2025-08-16, and the scripts are untracked files in it.
We therefore identify them by their SHA-256 hashes instead of a commit:
the hash of every copy is listed in \texttt{archived\_scripts.sha256}
of the release, and the table gives the hash of that list, of the DAX script
and of the code of the re-run.

\begin{center}\small
\begin{tabular}{p{0.58\textwidth}p{0.32\textwidth}}
\toprule
Artefact & SHA-256 (first 16 characters) \\
\midrule
\texttt{archived\_scripts.sha256 (all 122 archived scripts)} & \texttt{9a9820f888f3506a} \\
\texttt{D:\textbackslash CatBoost\textbackslash catBoost\_conGAN.py} & \texttt{650aa836792450c1} \\
\texttt{rerun\_leakfree.py} & \texttt{e7e3ef6369f73bbd} \\
\texttt{finish\_run.py} & \texttt{b18df2d9be00ad7f} \\
\texttt{leakfree/\_\_init\_\_.py} & \texttt{61852af947fe4ac4} \\
\texttt{leakfree/config.py} & \texttt{33a0f8c54386bb9f} \\
\texttt{leakfree/data.py} & \texttt{ea2fc646a53e10e3} \\
\texttt{leakfree/db.py} & \texttt{293765cf6410d903} \\
\texttt{leakfree/dm.py} & \texttt{4eb4b3adbf70fdf5} \\
\texttt{leakfree/evaluation.py} & \texttt{71c0b9181dd981b9} \\
\texttt{leakfree/models.py} & \texttt{b27afb27d60babfa} \\
\texttt{leakfree/portfolio.py} & \texttt{384530e87ee0e85b} \\
\texttt{leakfree/r3m4.py} & \texttt{617aeae0ba43dd50} \\
\texttt{leakfree/report.py} & \texttt{6583db4a8a82e1f9} \\
\texttt{leakfree/splits.py} & \texttt{7f4496f8aadfcbde} \\
\bottomrule
\end{tabular}
\end{center}

\subpoint{1.5}{The results of the re-run}{I strongly recommend rerunning the experiments using strictly chronological splits and training-only preprocessing, with an untouched final test set.}
The table below is the headline result: the equal-weight portfolio over
the stocks of each cell, with 20~bps per unit of turnover and a
stationary-bootstrap 95\% interval on the geometric return. The
interval is computed on the daily portfolio return, so it accounts for
the serial dependence of the daily series.

\begin{center}\small
\begin{tabular}{lrlrrrrr}
\toprule
Regime & Window & Model & Stocks & Days & Return & 95\% interval & Sharpe \\
\midrule
crash & 5 & CatBoost & 39 & 30 & -17.42\% & [-41.68, 14.99]\% & -2.70 \\
crash & 5 & LightGBM & 39 & 30 & -12.44\% & [-33.89, 12.22]\% & -2.19 \\
crash & 5 & RandomWalk & 39 & 30 & 0.00\% & [0.00, 0.00]\% & 0.00 \\
crash & 5 & XGBoost & 39 & 30 & -12.65\% & [-33.16, 11.37]\% & -2.30 \\
crash & 5 & cGAN & 39 & 30 & -12.80\% & [-28.13, 6.66]\% & -2.56 \\
\addlinespace
crash & 10 & CatBoost & 39 & 30 & -19.16\% & [-41.67, 9.55]\% & -2.92 \\
crash & 10 & LightGBM & 39 & 30 & -16.82\% & [-39.03, 9.71]\% & -2.75 \\
crash & 10 & RandomWalk & 39 & 30 & 0.00\% & [0.00, 0.00]\% & 0.00 \\
crash & 10 & XGBoost & 39 & 30 & -15.08\% & [-34.60, 8.61]\% & -2.51 \\
crash & 10 & cGAN & 39 & 30 & -17.49\% & [-36.33, 5.32]\% & -2.70 \\
\addlinespace
crash & 30 & CatBoost & 39 & 30 & -22.40\% & [-44.90, 6.50]\% & -2.98 \\
crash & 30 & LightGBM & 39 & 30 & -18.09\% & [-38.76, 7.55]\% & -2.85 \\
crash & 30 & RandomWalk & 39 & 30 & 0.00\% & [0.00, 0.00]\% & 0.00 \\
crash & 30 & XGBoost & 39 & 30 & -12.76\% & [-31.53, 10.80]\% & -2.15 \\
crash & 30 & cGAN & 39 & 30 & -16.43\% & [-37.84, 9.62]\% & -2.57 \\
\addlinespace
normal & 5 & CatBoost & 39 & 100 & -4.92\% & [-12.15, 2.88]\% & -1.95 \\
normal & 5 & LightGBM & 39 & 100 & -12.15\% & [-19.07, -4.50]\% & -5.91 \\
normal & 5 & RandomWalk & 39 & 100 & 0.00\% & [0.00, 0.00]\% & 0.00 \\
normal & 5 & XGBoost & 39 & 100 & -12.32\% & [-17.72, -7.45]\% & -7.02 \\
normal & 5 & cGAN & 39 & 100 & -16.98\% & [-21.00, -12.05]\% & -10.49 \\
\addlinespace
normal & 10 & CatBoost & 39 & 100 & -2.85\% & [-10.01, 4.00]\% & -1.00 \\
normal & 10 & LightGBM & 39 & 100 & -8.34\% & [-15.18, -1.63]\% & -3.98 \\
normal & 10 & RandomWalk & 39 & 100 & 0.00\% & [0.00, 0.00]\% & 0.00 \\
normal & 10 & XGBoost & 39 & 100 & -6.07\% & [-11.67, -0.26]\% & -2.82 \\
normal & 10 & cGAN & 39 & 100 & -14.63\% & [-18.28, -10.52]\% & -8.55 \\
\addlinespace
normal & 30 & CatBoost & 39 & 100 & -1.44\% & [-9.03, 6.48]\% & -0.48 \\
normal & 30 & LightGBM & 39 & 100 & -5.87\% & [-13.85, 2.80]\% & -2.18 \\
normal & 30 & RandomWalk & 39 & 100 & 0.00\% & [0.00, 0.00]\% & 0.00 \\
normal & 30 & XGBoost & 39 & 100 & -5.48\% & [-11.48, 0.86]\% & -2.45 \\
normal & 30 & cGAN & 39 & 100 & -14.10\% & [-18.18, -9.77]\% & -8.55 \\
\bottomrule
\end{tabular}
\end{center}

\begin{center}\small
\begin{tabular}{lrlrrrrr}
\toprule
Regime & Window & Model & RMSE & MASE & Return & Sharpe & Hit rate \\
\midrule
crash & 5 & CatBoost & 8.375 & 3.922 & -17.34\% & -1.55 & 47.2\% \\
crash & 5 & LightGBM & 8.363 & 3.871 & -12.45\% & -1.12 & 47.9\% \\
crash & 5 & RandomWalk & 2.781 & 0.980 & 0.00\% & 0.00 & 0.0\% \\
crash & 5 & XGBoost & 8.395 & 4.040 & -11.44\% & -1.15 & 48.0\% \\
crash & 5 & cGAN & 11.676 & 5.157 & -12.58\% & -1.04 & 45.1\% \\
\addlinespace
crash & 10 & CatBoost & 9.291 & 4.233 & -19.44\% & -1.75 & 45.8\% \\
crash & 10 & LightGBM & 8.661 & 3.995 & -16.61\% & -1.52 & 47.0\% \\
crash & 10 & RandomWalk & 2.781 & 0.980 & 0.00\% & 0.00 & 0.0\% \\
crash & 10 & XGBoost & 9.232 & 4.358 & -13.15\% & -1.39 & 47.1\% \\
crash & 10 & cGAN & 12.135 & 5.280 & -16.17\% & -1.52 & 44.8\% \\
\addlinespace
crash & 30 & CatBoost & 11.041 & 4.906 & -22.19\% & -2.12 & 44.3\% \\
crash & 30 & LightGBM & 9.521 & 4.270 & -17.26\% & -1.67 & 46.6\% \\
crash & 30 & RandomWalk & 2.781 & 0.980 & 0.00\% & 0.00 & 0.0\% \\
crash & 30 & XGBoost & 10.340 & 4.691 & -9.96\% & -1.11 & 47.3\% \\
crash & 30 & cGAN & 12.198 & 5.545 & -15.73\% & -1.48 & 45.5\% \\
\addlinespace
normal & 5 & CatBoost & 1.865 & 2.238 & -5.41\% & -0.59 & 42.5\% \\
normal & 5 & LightGBM & 1.907 & 2.256 & -12.66\% & -1.35 & 40.2\% \\
normal & 5 & RandomWalk & 0.853 & 1.001 & 0.00\% & 0.00 & 0.0\% \\
normal & 5 & XGBoost & 2.081 & 2.376 & -12.73\% & -1.40 & 40.7\% \\
normal & 5 & cGAN & 7.720 & 9.761 & -17.07\% & -1.94 & 39.8\% \\
\addlinespace
normal & 10 & CatBoost & 2.066 & 2.430 & -3.41\% & -0.34 & 42.5\% \\
normal & 10 & LightGBM & 1.931 & 2.268 & -8.95\% & -0.84 & 41.4\% \\
normal & 10 & RandomWalk & 0.853 & 1.001 & 0.00\% & 0.00 & 0.0\% \\
normal & 10 & XGBoost & 2.222 & 2.495 & -6.59\% & -0.71 & 42.1\% \\
normal & 10 & cGAN & 7.563 & 9.938 & -14.63\% & -1.96 & 39.2\% \\
\addlinespace
normal & 30 & CatBoost & 2.508 & 2.930 & -1.60\% & -0.09 & 42.9\% \\
normal & 30 & LightGBM & 2.072 & 2.420 & -6.17\% & -0.68 & 41.1\% \\
normal & 30 & RandomWalk & 0.853 & 1.001 & 0.00\% & 0.00 & 0.0\% \\
normal & 30 & XGBoost & 2.438 & 2.769 & -5.80\% & -0.67 & 41.5\% \\
normal & 30 & cGAN & 7.813 & 10.871 & -14.39\% & -1.71 & 40.4\% \\
\bottomrule
\end{tabular}
\end{center}

Returns are net of 20~bps per unit of turnover, averaged over the stocks
of the cell. MASE is the mean absolute error over the mean absolute
one-step change of the price, so a value above one means the model is
worse than a random walk on the same dates.

\subpoint{1.6}{The replication package}{Please provide an auditable replication package.}
The release contains the training program, the corrected evaluation
engine, the stored forecasts for every stock and model, the chronology of
every cell, and a hash of every file:

\begin{center}\small
\begin{tabular}{p{0.30\textwidth}p{0.58\textwidth}}
\toprule
Artefact & Content \\
\midrule
\texttt{rerun\_leakfree.py} & the re-run, with the split, the models and the storage \\
\texttt{leakfree/splits.py} & the chronological split and the train-only scaler \\
\texttt{leakfree/dm.py} & the Diebold--Mariano test in the three conventions in use \\
\texttt{leakfree/evaluation.py} & the corrected backtest and DM engine (Stage~1) \\
\texttt{data/chronology.csv} & train, validation and test dates per stock and cell \\
\texttt{data/predictions/} & the forecasts of every model, per stock \\
\texttt{data/metrics\_per\_stock.csv} & the backtest metrics of every model \\
\texttt{report/provenance.json} & hashes of inputs, code and libraries \\
\texttt{manifest.sha256} & the hash of every file in the release \\
\bottomrule
\end{tabular}
\end{center}

\section*{2. Response to Reviewer 3, comment R3.M3}

\point{2}{Market-regime definitions and evaluation periods}{R3.M3}
\begin{reviewercomment}
The market-regime definitions and evaluation periods are inconsistent. Section 3.10 describes a 100-day out-of-sample crash evaluation running from 17 February to 30 March 2020, although this calendar interval cannot contain 100 trading days. Section 4.1 subsequently reports only 31 common crash-period target dates. The normal-regime audit also contains uneven coverage: the manuscript reports partial W5 and W30 groups and no available normal W10 group. These differences make comparisons across algorithms, lookback windows, and regimes difficult to interpret.

In addition, the linked repository's README describes some forecasts labeled as \emph{normal} that cover crash-period dates, including identical stored files across the two labels for certain settings. The authors should clarify whether these are connected to the submitted results. Please provide a stock-by-stock and model-by-model chronology specifying training, validation, and test dates; the precise operational definition of each regime; the numbers of matched forecast observations; and any exclusions. All comparisons supporting regime-specific superiority should be recomputed on genuinely comparable dates and verified regime labels.
\end{reviewercomment}
\ourans
The reviewer is right that the regime definitions and the evaluation
periods were not stated consistently, and the corrected run fixes the
underlying problem rather than the description.

\subpoint{2.1}{The operational definition of each regime}{Section 3.10 describes a 100-day out-of-sample crash evaluation running from 17 February to 30 March 2020, although this calendar interval cannot contain 100 trading days, and Section 4.1 reports only 31 common crash-period target dates. The normal-regime audit also contains uneven coverage: partial W5 and W30 groups and no available normal W10 group. Please give the precise operational definition of each regime.}
Both regimes are defined by the forecast target dates, not by an index
offset, and the definition is the one of the manuscript's own group table:

\begin{center}\small
\begin{tabular}{p{0.16\textwidth}p{0.30\textwidth}p{0.44\textwidth}}
\toprule
Regime & Target dates & Definition and source \\
\midrule
Crash & 2020-02-17 to 2020-03-30 & the COVID-19 fall; group table G01--G15 of the supplementary material, which lists $T = 31$ forecast observations \\
Normal & 2023-11-08 to 2024-04-03 & the calm market after the COVID-19 episode; the manuscript's normal groups are partial ($T = 37$ at $W = 5$, $T = 12$ at $W = 30$) and the re-run uses the largest block every window and stock can support \\
\bottomrule
\end{tabular}
\end{center}

Section~3.10 of the manuscript describes the crash evaluation as a
100-day evaluation running from 17~February to 30~March 2020. That
interval contains 31 trading days, and the group table lists $T = 31$, so
the two statements cannot both hold. The revised manuscript now states the
dates and the number of observations and drops the 100-day description.

\subpoint{2.2}{The chronology, stock by stock and model by model}{Please provide a stock-by-stock and model-by-model chronology specifying training, validation, and test dates, the numbers of matched forecast observations, and any exclusions.}
Every cell of the re-run records the dates of its training, validation and
test blocks. The table in Section~0 gives the common part; the complete
stock-by-stock and model-by-model chronology, including the number of
matched forecast observations and every exclusion, is in
\texttt{data/chronology.csv} and \texttt{data/coverage.csv} of the release,
and in the audit website.

\subpoint{2.3}{Recomputation on genuinely comparable dates}{The linked repository's README describes some forecasts labeled as \emph{normal} that cover crash-period dates, including identical stored files across the two labels for certain settings. Please clarify whether these are connected to the submitted results. All comparisons supporting regime-specific superiority should be recomputed on genuinely comparable dates and verified regime labels.}
The corrected run evaluates every window of a regime on the same target
dates, so the cross-window comparisons in the revised manuscript rest on
matched observations. The archived runs did not: the crash and the normal
blocks at $W = 5$ and $W = 10$ were byte-identical prediction files, and
the normal block was truncated to the last rows of the sample. The
revised manuscript states the coverage of every cell and no longer
compares cells with different observation counts.

\section*{3. Response to Reviewer 3, comment R3.M4: the Diebold--Mariano statistics}

\point{3}{Definition and interpretation of the Diebold--Mariano statistics}{R3.M4}
\begin{reviewercomment}
Diebold--Mariano tests: clarification of the loss differential and interpretation. Section 3.7 defines the loss differential as
\[
d_t = g(e_{1,t}) - g(e_{2,t}),
\]
whereas Table 2 reports mostly positive Diebold--Mariano statistics for cGAN versus boosting models in the crash regime, which the discussion treats as evidence favoring the cGAN. Please make explicit which model corresponds to the first and second error series in every comparison, which loss function is used in the actual calculations, and how the sign of each statistic should be interpreted. For example, if the first error is the cGAN error, a positive mean squared-loss differential would conventionally indicate larger cGAN forecast errors. A worked example connecting one cell of Table 2 to the underlying paired losses would be helpful, especially in light of the RMSE rankings in Table 1.

Please also clarify how tests involving forecasts labeled as different regimes can satisfy the requirement that the competing models forecast the same realized outcomes on the same dates. The meaning of the reported cross-stock averages of test statistics, the treatment of serial and cross-sectional dependence, the small-sample correction, and any adjustment for multiple comparisons should be explained. I encourage the authors to report paired sample sizes, the mean loss differentials, uncertainty estimates, and adjusted $p$-values in addition to the test statistics. These requests are intended to make the statistical evidence transparent and interpretable.
\end{reviewercomment}
\ourans
While preparing the re-run we reverse engineered the published statistics
and can state exactly how they were produced, which the reviewer asks for
in R3.M4.

\subpoint{3.1}{The loss function and the published statistic}{Section 3.7 defines the loss differential as $d_t = g(e_{1,t}) - g(e_{2,t})$. Which loss function is used in the actual calculations?}
The published summary matrix
(\texttt{C:\textbackslash program\_stock\textbackslash dmtest\_summary\_matrix.xlsx}) was rebuilt
from the archived per-stock files and reproduces exactly: the largest
absolute difference over the nine cells of the first stock is
$4.4 \times 10^{-16}$, i.e.\ floating point noise. The published statistic
is
$$
d_t = e_{1,t}^2 - e_{2,t}^2, \qquad
z = \frac{\overline{d}}{\sqrt{\widehat{\mathrm{Var}}(d)/T}}, \qquad
\widehat{\mathrm{Var}}(d) = \frac{1}{T}\sum_t (d_t - \overline{d})^2,
$$
that is, the population variance of the loss differential and a standard
normal reference distribution, with no small-sample correction. The
MATLAB function \texttt{dmtest.m} in the repository uses the $N-1$
divisor instead, so its statistic is $\sqrt{(T-1)/T}$ times the published
one.

\subpoint{3.2}{Which model is first, and the sign of the statistic}{Please make explicit which model corresponds to the first and second error series in every comparison, and how the sign of each statistic should be interpreted. For example, if the first error is the cGAN error, a positive mean squared-loss differential would conventionally indicate larger cGAN forecast errors.}
The sign convention is stated in every table: $d_t = \mathrm{loss}_A - \mathrm{loss}_B$,
so a positive statistic means the first model has the larger loss and is
the worse one.

\subpoint{3.3}{Serial dependence, small-sample correction and multiple testing}{The treatment of serial and cross-sectional dependence, the small-sample correction, and any adjustment for multiple comparisons should be explained.}
The revision reports three statistics side by side: the published
convention, the MATLAB convention, and the corrected
Harvey--Leybourne--Newbold statistic with a Newey--West Bartlett HAC variance and a $t(T-1)$
reference distribution, with Holm adjustment inside each family.
For a one-step forecast the Newey--West lag of the per-stock test is
$h - 1 = 0$, so serial correlation of the loss differential is not
corrected there. To check this we repeated every per-stock test with
$\lfloor 4 (T/100)^{2/9} \rfloor$ Bartlett lags (3 for $T = 31$, 4 for $T = 101$). Of the 2,340 per-stock comparisons between two different models, 1,423 are significant after the Holm adjustment with lag~0 and 929 with the larger lag, so a part of the per-stock evidence rests on serially correlated differentials, and the lag-robust count is the more reliable of the two. Both counts are given for every cell in \nolinkurl{data/r3m4/pooled_tests.csv}.
The Holm adjustment controls the family-wise error rate under any
dependence between the tests of a family. Cross-sectional dependence
between stocks is treated in Point~3.6.

\subpoint{3.4}{A worked example}{A worked example connecting one cell of Table 2 to the underlying paired losses would be helpful, especially in light of the RMSE rankings in Table 1.}
We take ADVANC, the crash regime and $W = 10$, and compare cGAN (model A) with CatBoost (model B). Both models forecast the same 31 target dates, 2020-02-17 to 2020-03-30, so the realized prices are the same. With squared price errors the mean losses are 554.51 for cGAN and 76.23 for CatBoost (RMSE 23.548 and 8.731 THB), and the mean loss differential is $\bar d = 478.29$. The population variance of $d_t$ is 313,522.04, which gives the published statistic $z = \bar d / \sqrt{\widehat{\mathrm{Var}}(d)/T} = 4.756$. With the sample variance (323,972.78) the statistic is 4.679; for $h = 1$ this is also the Harvey--Leybourne--Newbold statistic, with $p = 5.8\times 10^{-5}$ against $t(30)$. The Holm adjustment over the 39 stocks of the cell gives $p = 0.0014$. With 3 Bartlett lags the statistic is 3.453 ($p = 0.0017$).
The positive sign means that cGAN has the larger loss: cGAN is the worse forecaster for this stock and cell.

\begin{center}\small
\begin{tabular}{lrrrrrr}
\toprule
Date & Price & cGAN & CatBoost & $e_{A,t}^2$ & $e_{B,t}^2$ & $d_t$ \\
\midrule
2020-02-17 & 214.00 & 170.43 & 199.93 & 1,898.00 & 198.08 & 1,699.92 \\
2020-02-18 & 213.00 & 220.43 & 202.54 & 55.14 & 109.40 & -54.26 \\
2020-02-19 & 210.00 & 215.16 & 203.43 & 26.62 & 43.18 & -16.56 \\
2020-02-20 & 206.00 & 209.21 & 203.35 & 10.30 & 7.00 & 3.30 \\
2020-02-21 & 211.00 & 231.80 & 203.57 & 432.75 & 55.25 & 377.50 \\
\bottomrule
\end{tabular}
\end{center}

The table shows the first five of the 31 paired dates; all of them are in
\nolinkurl{data/r3m4/worked_example_paired_losses.csv}, and the statistics are
the row of \nolinkurl{data/dm_per_stock.csv} for this stock and pair.
Because $\bar d$ is the difference of the two mean squared errors on the same
dates, the sign of the statistic always agrees with the RMSE ranking on those
dates. A statistic whose sign disagrees with an RMSE ranking, as the reviewer
observes for Tables~1 and~2 of the submitted manuscript, means that the two
were computed on different dates or with different losses.

\subpoint{3.5}{Forecasts of different windows and regimes}{Please also clarify how tests involving forecasts labeled as different regimes can satisfy the requirement that the competing models forecast the same realized outcomes on the same dates.}
In the published summary matrix behind Table~2
(\texttt{dmtest\_summary\_matrix.xlsx}) each cell compared CatBoost at one
window with LightGBM at one of the three windows, and the two error series
were paired by row position (\texttt{leakfree/dm.py},
\texttt{paper\_matrix\_cell\_errors}, which reproduces the published matrix
exactly). The archived scripts locate the
test block by an index offset into the windows (Point~1.1), and the number of
windows depends on the window length, so the test blocks of the three windows
start on different dates. The off-diagonal cells therefore paired errors of
different dates, and the requirement the reviewer states was not met.

In the re-run every test pairs two models of the same stock, the same regime
and the same window. The two forecasts are columns of the same prediction file,
which has one column of realized prices, so they refer to the same realized
outcome on the same date by construction, and no test compares forecasts of
different regimes. Within a regime the three windows have the same target
dates (the table in Section~0), so a comparison across windows is also made on
the same dates.

\subpoint{3.6}{Cross-stock averages and cross-sectional dependence}{The meaning of the reported cross-stock averages of test statistics should be explained.}
The submitted Table~2 reported, for each cell, the average of the per-stock
statistics over the stocks. Such an average is not a test statistic. Under the
null hypothesis each per-stock statistic is approximately standard normal, but
the average of $N$ of them has variance $N^{-2}\sum_{i,j}\rho_{ij}$, which is
$1/N$ only if the stocks are independent and tends to the average correlation
$\bar\rho$ when they are not. The prices of SET50 stocks move together,
particularly in the crash, so the average cannot be compared with a normal
critical value. The revision keeps the average only as a descriptive column.

Inference rests on two tests instead. The first is the per-stock test with the
Holm adjustment, which is valid under any dependence between stocks. The
second is a pooled test that allows for cross-sectional dependence: on each
date the loss differential is averaged over the stocks of the cell, and the
test is applied to that daily average with a Newey--West variance with
$\lfloor 4 (T/100)^{2/9} \rfloor$ Bartlett lags, the Harvey--Leybourne--Newbold
factor and a $t(T-1)$ reference distribution. Because the average is formed
date by date before the variance is estimated, the covariances between the
stocks enter the variance. The errors are scaled by the previous close, so the
pooled loss is a squared return error in percent squared and stocks with high
prices do not dominate the average. The $p$-values of the pooled test are
Holm-adjusted over the model pairs of each cell.

\subpoint{3.7}{Sample sizes, mean loss differentials, intervals and adjusted $p$-values}{I encourage the authors to report paired sample sizes, the mean loss differentials, uncertainty estimates, and adjusted $p$-values in addition to the test statistics.}
The table reports, for every cell and every pair of models, the number of
stocks $N$, the number of paired target dates $T$, the mean loss
differential $\bar d$ (squared return error, percent squared) with its 95\%
interval, the pooled statistic with its Holm-adjusted $p$-value, the average
of the per-stock statistics (descriptive only, Point~3.6) and the number of
stocks whose difference is significant after the Holm adjustment, split by
the model with the larger loss. A positive $\bar d$ means that model A has the
larger loss. CB is CatBoost, LGB LightGBM, XGB XGBoost and RW the random
walk.

{\footnotesize\setlength{\tabcolsep}{4pt}
\begin{longtable}{lrlrrrrrrrr}
\toprule
Regime & $W$ & A -- B & $N$ & $T$ & $\bar d$ & 95\% interval & Pooled & $p_{\mathrm{Holm}}$ & Avg.\ $z$ & A / B worse \\
\midrule
\endhead
\bottomrule
\endfoot
crash & 5 & cGAN -- CB & 39 & 31 & 101.6 & [38.97, 164.3] & 3.26 & 0.022 & 0.22 & 17 / 9 \\
crash & 5 & cGAN -- LGB & 39 & 31 & 86.18 & [19.89, 152.5] & 2.61 & 0.042 & 0.15 & 18 / 9 \\
crash & 5 & cGAN -- XGB & 39 & 31 & -42.79 & [-139.6, 54.01] & -0.89 & 0.382 & 0.09 & 17 / 9 \\
crash & 5 & cGAN -- RW & 39 & 31 & 1,119.1 & [430.2, 1,808.0] & 3.26 & 0.022 & 4.62 & 36 / 0 \\
crash & 5 & CB -- LGB & 39 & 31 & -15.47 & [-37.47, 6.52] & -1.41 & 0.336 & -0.56 & 6 / 10 \\
crash & 5 & CB -- XGB & 39 & 31 & -144.4 & [-227.3, -61.57] & -3.50 & 0.013 & -2.36 & 3 / 13 \\
crash & 5 & CB -- RW & 39 & 31 & 1,017.5 & [339.8, 1,695.2] & 3.02 & 0.027 & 3.70 & 17 / 0 \\
crash & 5 & LGB -- XGB & 39 & 31 & -129.0 & [-194.3, -63.68] & -3.97 & 0.004 & -2.44 & 4 / 12 \\
crash & 5 & LGB -- RW & 39 & 31 & 1,032.9 & [337.6, 1,728.3] & 2.98 & 0.027 & 3.28 & 16 / 0 \\
crash & 5 & XGB -- RW & 39 & 31 & 1,161.9 & [402.0, 1,921.8] & 3.07 & 0.027 & 3.41 & 17 / 0 \\
\addlinespace
crash & 10 & cGAN -- CB & 39 & 31 & 56.88 & [25.93, 87.84] & 3.69 & 0.008 & -0.20 & 18 / 12 \\
crash & 10 & cGAN -- LGB & 39 & 31 & 89.43 & [50.76, 128.1] & 4.65 & $<0.001$ & -0.06 & 19 / 12 \\
crash & 10 & cGAN -- XGB & 39 & 31 & -145.9 & [-242.2, -49.59] & -3.04 & 0.025 & -0.24 & 19 / 12 \\
crash & 10 & cGAN -- RW & 39 & 31 & 1,155.3 & [404.0, 1,906.6] & 3.09 & 0.025 & 4.96 & 36 / 0 \\
crash & 10 & CB -- LGB & 39 & 31 & 32.55 & [6.46, 58.63] & 2.51 & 0.025 & 2.07 & 13 / 4 \\
crash & 10 & CB -- XGB & 39 & 31 & -202.8 & [-319.7, -85.89] & -3.49 & 0.011 & -1.99 & 8 / 13 \\
crash & 10 & CB -- RW & 39 & 31 & 1,098.4 & [367.6, 1,829.3] & 3.02 & 0.025 & 4.07 & 22 / 0 \\
crash & 10 & LGB -- XGB & 39 & 31 & -235.3 & [-365.5, -105.2] & -3.63 & 0.008 & -3.25 & 1 / 15 \\
crash & 10 & LGB -- RW & 39 & 31 & 1,065.9 & [351.0, 1,780.8] & 3.00 & 0.025 & 3.47 & 17 / 0 \\
crash & 10 & XGB -- RW & 39 & 31 & 1,301.2 & [456.8, 2,145.7] & 3.10 & 0.025 & 3.62 & 18 / 0 \\
\addlinespace
crash & 30 & cGAN -- CB & 39 & 31 & 52.77 & [-28.77, 134.3] & 1.30 & 0.203 & -0.43 & 11 / 9 \\
crash & 30 & cGAN -- LGB & 39 & 31 & 183.5 & [55.24, 311.7] & 2.87 & 0.037 & 0.22 & 18 / 9 \\
crash & 30 & cGAN -- XGB & 39 & 31 & -98.03 & [-174.6, -21.50] & -2.57 & 0.037 & -0.07 & 20 / 13 \\
crash & 30 & cGAN -- RW & 39 & 31 & 1,306.5 & [433.6, 2,179.4] & 3.01 & 0.037 & 5.14 & 39 / 0 \\
crash & 30 & CB -- LGB & 39 & 31 & 130.7 & [40.50, 220.9] & 2.91 & 0.037 & 1.94 & 24 / 5 \\
crash & 30 & CB -- XGB & 39 & 31 & -150.8 & [-235.5, -66.11] & -3.58 & 0.012 & -0.25 & 16 / 11 \\
crash & 30 & CB -- RW & 39 & 31 & 1,253.7 & [409.0, 2,098.5] & 2.98 & 0.037 & 4.84 & 30 / 0 \\
crash & 30 & LGB -- XGB & 39 & 31 & -281.5 & [-447.8, -115.2] & -3.40 & 0.017 & -3.78 & 3 / 19 \\
crash & 30 & LGB -- RW & 39 & 31 & 1,123.0 & [365.4, 1,880.7] & 2.98 & 0.037 & 3.80 & 22 / 0 \\
crash & 30 & XGB -- RW & 39 & 31 & 1,404.5 & [481.0, 2,328.1] & 3.06 & 0.037 & 3.96 & 21 / 0 \\
\addlinespace
normal & 5 & cGAN -- CB & 39 & 101 & 195.1 & [171.2, 218.9] & 16.15 & $<0.001$ & 4.93 & 33 / 4 \\
normal & 5 & cGAN -- LGB & 39 & 101 & 194.0 & [170.1, 217.8] & 16.06 & $<0.001$ & 4.86 & 33 / 4 \\
normal & 5 & cGAN -- XGB & 39 & 101 & 191.7 & [167.8, 215.7] & 15.83 & $<0.001$ & 4.61 & 33 / 4 \\
normal & 5 & cGAN -- RW & 39 & 101 & 219.4 & [195.2, 243.6] & 17.87 & $<0.001$ & 7.03 & 37 / 0 \\
normal & 5 & CB -- LGB & 39 & 101 & -1.10 & [-1.97, -0.24] & -2.53 & 0.013 & -0.60 & 4 / 5 \\
normal & 5 & CB -- XGB & 39 & 101 & -3.33 & [-4.67, -1.99] & -4.90 & $<0.001$ & -1.39 & 3 / 7 \\
normal & 5 & CB -- RW & 39 & 101 & 24.32 & [15.42, 33.21] & 5.39 & $<0.001$ & 3.70 & 14 / 0 \\
normal & 5 & LGB -- XGB & 39 & 101 & -2.22 & [-2.89, -1.55] & -6.55 & $<0.001$ & -1.16 & 4 / 5 \\
normal & 5 & LGB -- RW & 39 & 101 & 25.42 & [15.85, 34.99] & 5.24 & $<0.001$ & 3.80 & 19 / 0 \\
normal & 5 & XGB -- RW & 39 & 101 & 27.64 & [17.69, 37.59] & 5.48 & $<0.001$ & 3.69 & 14 / 0 \\
\addlinespace
normal & 10 & cGAN -- CB & 39 & 101 & 228.2 & [201.8, 254.6] & 17.07 & $<0.001$ & 5.44 & 34 / 2 \\
normal & 10 & cGAN -- LGB & 39 & 101 & 230.8 & [204.3, 257.3] & 17.20 & $<0.001$ & 5.61 & 34 / 2 \\
normal & 10 & cGAN -- XGB & 39 & 101 & 225.5 & [199.4, 251.7] & 17.01 & $<0.001$ & 5.38 & 34 / 2 \\
normal & 10 & cGAN -- RW & 39 & 101 & 256.2 & [225.3, 287.2] & 16.35 & $<0.001$ & 6.81 & 38 / 0 \\
normal & 10 & CB -- LGB & 39 & 101 & 2.59 & [1.81, 3.37] & 6.53 & $<0.001$ & 3.02 & 11 / 2 \\
normal & 10 & CB -- XGB & 39 & 101 & -2.66 & [-4.20, -1.12] & -3.41 & $<0.001$ & -0.39 & 6 / 5 \\
normal & 10 & CB -- RW & 39 & 101 & 28.03 & [18.29, 37.78] & 5.68 & $<0.001$ & 4.18 & 17 / 0 \\
normal & 10 & LGB -- XGB & 39 & 101 & -5.25 & [-6.61, -3.88] & -7.58 & $<0.001$ & -1.58 & 2 / 5 \\
normal & 10 & LGB -- RW & 39 & 101 & 25.44 & [15.84, 35.05] & 5.23 & $<0.001$ & 3.70 & 17 / 0 \\
normal & 10 & XGB -- RW & 39 & 101 & 30.69 & [19.90, 41.49] & 5.61 & $<0.001$ & 3.76 & 15 / 0 \\
\addlinespace
normal & 30 & cGAN -- CB & 39 & 101 & 331.3 & [284.3, 378.2] & 13.93 & $<0.001$ & 5.83 & 34 / 2 \\
normal & 30 & cGAN -- LGB & 39 & 101 & 340.5 & [293.4, 387.5] & 14.28 & $<0.001$ & 6.18 & 34 / 2 \\
normal & 30 & cGAN -- XGB & 39 & 101 & 330.5 & [283.7, 377.2] & 13.96 & $<0.001$ & 5.42 & 34 / 3 \\
normal & 30 & cGAN -- RW & 39 & 101 & 370.3 & [320.4, 420.2] & 14.65 & $<0.001$ & 7.49 & 39 / 0 \\
normal & 30 & CB -- LGB & 39 & 101 & 9.20 & [7.01, 11.40] & 8.27 & $<0.001$ & 2.25 & 14 / 4 \\
normal & 30 & CB -- XGB & 39 & 101 & -0.80 & [-3.70, 2.11] & -0.54 & 0.589 & 0.10 & 10 / 5 \\
normal & 30 & CB -- RW & 39 & 101 & 39.06 & [27.12, 51.00] & 6.46 & $<0.001$ & 5.05 & 24 / 0 \\
normal & 30 & LGB -- XGB & 39 & 101 & -10.00 & [-12.80, -7.20] & -7.05 & $<0.001$ & -2.82 & 2 / 8 \\
normal & 30 & LGB -- RW & 39 & 101 & 29.85 & [19.13, 40.57] & 5.50 & $<0.001$ & 3.97 & 19 / 0 \\
normal & 30 & XGB -- RW & 39 & 101 & 39.85 & [26.40, 53.30] & 5.85 & $<0.001$ & 4.26 & 18 / 0 \\
\end{longtable}}

For the cGAN the table reads as follows, with the Holm-adjusted pooled $p$-value at the 5\% level. In the crash regime the pooled test gives the cGAN a significantly smaller loss than XGBoost at $W = 10$ and $W = 30$; a significantly larger loss than CatBoost, LightGBM and the random walk at $W = 5$ and $W = 10$, and LightGBM and the random walk at $W = 30$ (2 of 12 comparisons in its favour, 8 against it, 2 not significant). In the normal regime the pooled test gives the cGAN a significantly larger loss than every other model at every window (0 of 12 comparisons in its favour, 12 against it, 0 not significant).

The table is \nolinkurl{data/r3m4/pooled_tests.csv} of the release, which
also gives the standard errors, the unadjusted $p$-values, the number of
paired observations and the per-stock counts with the lag-robust variance.

\end{document}